\lecture{6}
\title{CT Fourier Transform}

\begin{document}

\subsection{Lecture: Motivating the CT Fourier Transform}

Our goal is to write Fourier series for \emph{aperiodic} signals.
Here's our prototypical example.
\[ f(t) := \begin{cases}
1 & \text{ if }  t \in (-1, 1). \\
0 & \text{ otherwise.}
\end{cases} \]
The guiding intuition is that all signals are technically periodic—it's just that
the period of an aperiodic signal is $T \to \infty$.
Instead of a \emph{finite sum} $\sum_{k = -\infty}^{\infty} F[k] e^{jkt \omega_0}$,
we'll write $f$ as an \emph{infinite integral}
$\int_{-\infty}^{\infty} F(\omega) e^{jt \omega} \, \mathrm{d}\omega$.

But let's take things slow.
To build up to being able to say that $f(t)$ has period $T \to \infty$,
let's start by considering what happens when we force $f(t)$
to have some \emph{finite} period $T$ by constructing an alternative $f_T(t)$.
\[ f_T(t) := \begin{cases}
1 & \text{ if } (t \% T) \in (-1, 1). \\
0 & \text{ otherwise.}
\end{cases} \]
As a function of $T$, we now determine the length-$T$ array of Fourier coefficients $F_T$
for the periodic signal $f_T(t)$.
\[ F_T[k] = \frac{1}{T} \int_{-1}^1 e^{-jt(\textcolor{purple}{2\pi k/T})} \, \mathrm{d}t
= \frac{1}{T} \left[
    \frac{e^{-jt(\textcolor{purple}{2\pi k/T})}}{-j(\textcolor{purple}{2\pi k/T})}
\right]_{-1}^{1}
= \frac{1}{T} \cdot \frac{2\sin\left(\textcolor{purple}{2\pi k/T}\right)}{(\textcolor{purple}{2\pi k/T})}. \]
And here's a plot of $T \cdot F_T[k] =
\frac{2\sin\left(\textcolor{purple}{2\pi k/T}\right)}{(\textcolor{purple}{2\pi k/T})}$
for $T = 16$ and $T = 8$.
(Note that $k$ can be \emph{any} integer, not just those pictured below.)
\begin{center}
    \frame{\includegraphics[width=0.45\textwidth]{lecture-06/coefficients}}
\end{center}

\begin{remark}
    Convince yourself that doubling $T$ has the effect of
    halving the ``height'' and doubling the ``width''.
    If we want to ``normalize'' the graphs for $T = 16$ and $T = 8$ so they look similar,
    we'd instead have to graph, say, $T \cdot F_T[k / T]$\dots
\end{remark}

And so our Fourier series for $f_T(t)$ looks like the following:
\[ f_T(t) = \frac{1}{T} \sum_{k = -\infty}^{\infty}
\frac{2\sin(\textcolor{purple}{2\pi k/T})}{\textcolor{purple}{2\pi k/T}}
e^{jt(\textcolor{purple}{2\pi k/T})}. \]
There's a reparametrization begging to be made here:
$\textcolor{purple}{\{\omega_k\}_{-\infty}^{\infty} := 2\pi k/T}$.
Note that $\Delta \omega := \omega_{k + 1} - \omega_k = 2\pi/T$.
Now we can send $T \to \infty$ and see what happens to the Fourier series.
\begin{align*}
f(t) = \lim_{T \to \infty} f_T(t) & = \lim_{T \to \infty} \frac{1}{T} \sum_{k = -\infty}^{\infty}
\frac{2\sin(\textcolor{purple}{2\pi k/T})}{\textcolor{purple}{2\pi k/T}}
e^{jt(\textcolor{purple}{2\pi k/T})} \\
& = \lim_{T \to \infty} \frac{1}{2\pi} \sum_{k = -\infty}^{\infty}
\frac{2\sin(\textcolor{purple}{\omega_k})}{\textcolor{purple}{\omega_k}}
e^{jt \textcolor{purple}{\omega_k}} \Delta \omega
= \frac{1}{2\pi}\int_{-\infty}^{\infty}
\frac{2\sin(\omega)}{\omega}
e^{jt \omega} \, \mathrm{d}\omega.
\end{align*}
And so our Fourier series now becomes a Fourier integral!
\begin{center}
    \frame{\includegraphics[width=0.45\textwidth]{lecture-06/sum-to-integral}}
\end{center}

\subsection{Lecture: The CT Fourier Transform}

Naturally, we'd like to generalize:
what's the relationship between $f(t)$ and $F(\omega) = \frac{2\sin(\omega)}{\omega}$,
and vice versa?

\begin{theorem}{Fourier Transform}
    Given a CT signal $f(t)$, define its Fourier transform to be
    $F(\omega) := \int_{-\infty}^{\infty} f(t) e^{-j\omega t} \, \mathrm{d}t$.
    It then follows that
    $f(t) = \frac{1}{2\pi} \int_{-\infty}^{\infty} F(w) e^{jt \omega} \, \mathrm{d}\omega$.
\end{theorem}

\begin{proof}
    Follow the same recipe as before.
    Let $f_T(t)$ denote the extension of $f(t)$ to have period $T$.
    Then its Fourier series is:
    \[ F_T[k] = \frac{1}{T} \int_T f(t) e^{-jt(\textcolor{purple}{2\pi k/T})} \, \mathrm{d}t \]
    In our prototypical example,
    we picked $F(\omega) = \frac{2\sin(\omega)}{\omega}$ so that
    $F(\omega) := \lim_{T \to \infty} T \cdot F_T[k]$, where
    $\textcolor{purple}{\omega = 2\pi k/T}$.
    Let's follow the recipe and see what happens.
    \[ F(\omega) = \lim_{T \to \infty} T \cdot
    \left(\frac{1}{T} \int_T f(t) e^{-jt(\textcolor{purple}{2\pi k/T})} \, \mathrm{d}t\right)
    = \lim_{T \to \infty} \int_T f(t) e^{-jt\textcolor{purple}{\omega}} \, \mathrm{d}t
    = \int_{-\infty}^{\infty} f(t) e^{-jt\textcolor{purple}{\omega}} \, \mathrm{d}t. \]
    Just as promised.  With $F(\omega)$ defined this way,
    watch what happens when we take the Fourier series of $f(t)$ and send $T \to \infty$.
    \begin{align*}
        f(t) = \lim_{T \to \infty} f_T(t) & =
        \lim_{T \to \infty}
        \sum_{k = -\infty}^{\infty} F_T[k] e^{jt(\textcolor{purple}{2\pi k/T})} \\
        & = \frac{1}{2\pi} \lim_{T \to \infty} \sum_{k = -\infty}^{\infty} F(\omega_k)
        e^{jt(\textcolor{purple}{\omega_k})} \Delta \omega
        = \frac{1}{2\pi}\int_{-\infty}^{\infty} F(\omega) e^{jt\omega} \, \mathrm{d}\omega.
    \end{align*}
    Just as promised. ~ $\blacksquare$
\end{proof}

\begin{remark}
    Comparing and contrasting periodic and aperiodic signals\dots
    \begin{itemize}
        \item All the frequencies in a periodic signal
        are captured in one period $\left(\int_T\right)$,
        whereas aperiodic signals need all of $\mathbb{R}$
        $\left(\int_{-\infty}^{\infty}\right)$.
        \item Periodic signals have a discrete set of frequencies
        $\left(\sum_{k = -\infty}^{\infty}\right)$, whereas
        aperiodic signals have a continuous set
        $\left(\int_{-\infty}^{\infty} \, \mathrm{d}\omega\right)$.
    \end{itemize}
\end{remark}

Importantly, there is now a more clear \emph{dual} relationship
between $f(t)$ and $F(\omega)$.  If the equation
$F(\omega) := \int_{-\infty}^{\infty} f(t) e^{-j\omega t} \, \mathrm{d}t$
represents a Fourier transform, then the equation
$f(t) = \frac{1}{2\pi} \int_{-\infty}^{\infty} F(w) e^{jt \omega} \, \mathrm{d}\omega$
represents an inverse Fourier transform.

\begin{theorem}{Time-Frequency Duality}
    If the Fourier transform of $f(t)$ is $F(\omega)$,
    then the Fourier transform of $F(t)$ is $2\pi f(-\omega)$.
\end{theorem}

\begin{proof}
    Just look at the formulas. ~ $\blacksquare$
\end{proof}

\subsection{Lecture and Recitation: Properties of the CT Fourier Transform}

\begin{theorem}{Time Delay}
    The FT of $f(t - \tau)$ is $e^{-j \omega \tau} F(\omega)$,
    and the inverse FT of $F(\omega - \omega_0)$ is $e^{jt \omega_0}f(t)$.
\end{theorem}

\begin{proof}
    True by inspecting the relationship
    $F(\omega) := \int_{-\infty}^{\infty} f(t) e^{-j\omega t} \, \mathrm{d}t$. ~ $\blacksquare$
\end{proof}

Note that this suggests that our prototypical example
was unique in having a real-valued Fourier transform.  Here's why.

\begin{theorem}{Even / Odd}
    If $f(t)$ is real and even, $F(\omega)$ is real and even.
    If $f(t)$ is real and odd, $F(\omega)$ is imaginary and odd.
\end{theorem}

\begin{proof}
    Again, true by inspecting the relationship
    $F(\omega) := \int_{-\infty}^{\infty} f(t) e^{-j\omega t} \, \mathrm{d}t$.
    Another reason to see why is to apply the same argument as used for periodic signals;
    even signals are sums of pure cosines, and odd signals are sums of pure sines.
    ~ $\blacksquare$
\end{proof}

\begin{theorem}{Time Scaling}
    The FT of $f(at)$ is $\frac{1}{|a|}F\left(\frac{\omega}{a}\right)$,
    and the inverse FT of $F(a\omega)$ is $\frac{1}{|a|}f\left(\frac{1}{a}t\right)$.
\end{theorem}

\begin{proof}
    Again, just inspect the relationship
    $F(\omega) := \int_{-\infty}^{\infty} f(t) e^{-j\omega t} \, \mathrm{d}t$.
    The absolute value signs $|a|$ are necessary because
    when $a$ is negative, the $u$-substitution $u = at$ causes the bounds of the integral to flip.

    Alternatively, we can argue heuristically.
    This relationship says that, say,
    doubling the width of $f(t)$ causes its FT to
    shrink in width and double in height.
    \begin{center}
        \frame{\includegraphics[width=0.45\textwidth]{lecture-06/time-scaling}}
    \end{center}
    This should make sense; the signal $f(t/2)$ can be thought of as
    ``playing back the signal $f(t)$ at $\times 0.5$ speed'',
    so the frequencies prominent in $f(t/2)$ should be half of those prominent in $f(t)$.
    (And to balance it out, all frequency strengths are doubled.) ~ $\blacksquare$
\end{proof}

Here's some intuition as to why ``balancing it out'' is necessary.

\begin{theorem}{FT Areas}
    The area under the curve of $f(t)$ is $F(0)$,
    and the area under the curve of $F(\omega)$ is $2\pi f(0)$.
\end{theorem}

\begin{proof}
    Immediate from plugging in $t = 0$ or $\omega = 0$ as needed. ~ $\blacksquare$
\end{proof}

Finally, here's a property that \emph{doesn't}
have an analogue for Fourier series.

\begin{theorem}{Derivative FTs}
    The FT of $\frac{d}{dt} f(t)$ is $j\omega F(\omega)$,
    and the inverse FT of $\frac{d}{d\omega} F(\omega)$ is $-jt f(t)$.
\end{theorem}

\begin{proof}
    This one is trickier.  To show the FT of $\frac{d}{dt} f(t)$ is $j\omega F(\omega)$,
    differentiate the synthesis equation.
    \[ \frac{d}{dt} f(t)
    = \frac{d}{dt} \frac{1}{2\pi} \int_{-\infty}^{\infty} F(\omega) e^{jt\omega} \, \mathrm{d}\omega
    = \frac{1}{2\pi} \int_{-\infty}^{\infty} j\omega F(\omega) e^{jt\omega} \, \mathrm{d}\omega. \]
    To show inverse FT of $\frac{d}{d\omega} F(\omega)$ is $-jt f(t)$,
    differentiate the analysis equation.
    \[ \frac{d}{d\omega} F(\omega)
    = \frac{d}{d\omega} \int_{-\infty}^{\infty} f(t) e^{-jt \omega} \, \mathrm{d}t
    = \int_{-\infty}^{\infty} -jt f(t) e^{-jt \omega} \, \mathrm{d}t. ~~~ \blacksquare \]
\end{proof}

\subsection{Lecture: The Dirac Delta Function}

Consider our prototypical function $f(t)$ from earlier.
Imagine compressing $f(t)$ inward by a factor of $A$, producing functions like:
\[ f_A(t) := \begin{cases}
    A & \text{ if } t \in \left(-\frac{1}{A}, \frac{1}{A}\right) \\
    0 & \text{ otherwise.}
\end{cases} \]
Then the area under the curve of $f_A(t)$ is constantly $2$,
so $F_A(0)$ must be constantly $2$.
According to ``Time Scaling'', in fact,
the effect is that $F_A(\omega)$ ``stretches outward''
as $A$ tends to infinity, while still maintaining a $y$-intercept of $2$.

At the limit of $A \to \infty$, we find the following.

\begin{theorem}{Dirac Delta FT}
    Let the \textbf{Dirac delta function} $\delta(t)$
    be defined so that $\delta(t) = 0$ for all $t \neq 0$
    and $\int_{-\infty}^{\infty} \delta(t) \, \mathrm{d}t = 1$.
    Then the FT of $\delta(t)$ is the constant function $F(\omega) = 1$.

    The dual also holds: the FT of the constant function $f(t) = 1$
    is $F(\omega) = 2\pi \delta(\omega)$.
\end{theorem}

\begin{proof}
    This follows from the heuristic argument from earlier.
    You can also formally verify this using the FT equations. ~ $\blacksquare$
\end{proof}

The Dirac delta function also appears
when taking the Fourier transform of a periodic signal.

\begin{theorem}{Periodic FT}
    Suppose $f(t)$ is a periodic signal with Fourier coefficients $F[k]$,
    so that $f(t) = \sum_{k = -\infty}^{\infty} F[k] e^{jt(2\pi k/T)}$.
    Then the Fourier transform of $f(t)$ is
    $F(\omega) := \sum_{k = -\infty}^{\infty} 2\pi F[k] \delta\left(\omega - \frac{2\pi k}{T}\right)$.
\end{theorem}

\begin{proof}
    Nicely, the FT is linear, so it suffices to find the FT of $F[k] e^{jt(2\pi k/T)}$.
    This can be done nicely via ``Time Delay''.
    \begin{center}
        \frame{\includegraphics[width=0.45\textwidth]{lecture-06/periodic-FT}}
    \end{center}
    This should also just heuristically make sense:
    integrating over a discrete set of Dirac deltas
    is the same as summing over a discrete set of complex exponentials.
    ~ $\blacksquare$
\end{proof}

\end{document}